\section{边缘并不在地图外：辽东、高句丽与海峡}

\subsection{辽东的选择不是魏吴二选一}

公孙氏统治的辽东夹在中原道路、东北山地和海上消息之间。孙权给公孙渊的册命
与礼物能够抵达，说明建业愿意把这处边缘变成牵制魏的接口；它并不等于孙吴
已经拥有可以穿越海峡、长期保护盟友的军政能力。公孙渊衡量的也不仅是名号：
城池、粮食、部众、魏军的接近和海上援助的不确定性，都会使一次接受册拜成为
可被随时改写的风险计算。

吴使张弥、许晏遇害后，海上联络断裂。《三国志》及裴注所引材料能确认结果，
但不能让我们进入辽东每一座城或每一群人的意见。把事件写成公孙渊“反复无常”
很省力，却抹掉了远方政治的物质限制。册封的纸、印绶和礼物必须经船、港口、
翻译与接待者转送；任何一环失效，建业的承诺都无法变成襄平附近的保护。

司马懿围襄平时，魏的远征也不应被看成中原自然延伸到东北。辽河平原的雨水、
道路、围城粮食和回程，都在决定军队能停留多久。公孙渊的败亡结束了其政权的
核心，却没有把辽东居民、邻近山地群体和地方信息同时变成洛阳的财产。任官、
安抚、驻防和与高句丽的关系才是军事进入后的问题；中央战报对这些缓慢过程的
保存远少于对胜利的书写。

\subsection{丸都城之外，高句丽不是魏军报告的背景}

毌丘俭进攻高句丽、攻破丸都城，能够确认两个政治体在辽东山地发生过强制接触。
魏军需沿山路推进，高句丽的王和城众则在山城、河谷与补给条件中作出避走或
抵抗的选择。一次破城能够改变当年的军事秩序，却不能证明魏已长期控制当地
土地、人口或语言。魏军必须回转，便说明远征能力与持续统治之间存在一条不应
由胜利叙事抹去的距离。

后出的《三国史记》与区域考古为地名、山城和空间问题提供另一组线索，但其
成书层次不同，不能同魏方纪传无缝拼接为逐日行军图。考古中的城址、器物和
层位能够帮助理解山地聚落与防御的可能条件，却不会自动标明某一支部队的路线
或某位统治者的意图。保留两类材料的不重合，反而使丸都不再只是中原本纪里
一笔可略过的战功。

\subsection{带方与海峡：卑弥呼的使者为何重要}

卑弥呼遣使到魏，首先是一段必须经过带方、港口、译者和海峡的往返。魏廷授予
名号和物品，使远方来使被纳入朝贡礼仪；这种礼仪是双方都可利用的政治资源，
不是跨海行政的替代品。使者能否回返、地方中介怎样转译、岛屿一侧如何使用
这些物品，现存材料没有提供连续记录。故不能从《魏志·倭人传》的路线文字
推成稳定航线，也不能把后世国家边界投回三世纪海面。

这段接触仍值得与襄平围城并列。二三八年的魏一面以陆军处理辽东政权，一面
通过来使获得海上消息；两者显示的是不同的可达性。陆军可以围一座城，未必
能替海峡另一侧安排日常秩序；册封可以被带回远方，未必能让洛阳指挥船只和
港口。把这两种关系区别开来，才不会把东北写成一片由魏纪年自动填满的边疆。

\subsection{海上沉默与读者的责任}

孙吴的远航、魏与倭的往来，以及辽东与建业之间失败的外交，留下的多是朝廷
文件和胜负叙事。船工、沿海居民、译者和港口中介出现时，常只在他们使国家
计划得以或不能得以完成的瞬间。我们不能用这份沉默给他们编造统一意志，也不
应因此把他们从故事中删去。每次越海或沿江的行动，都需他们提供知识、劳力、
停泊与转送；这已经足以改变我们对“三国外交”的理解。

因此，图上的海峡应保留为关系而非边界。图件可以提示带方、辽东、建业和岛屿
之间的相对位置，不能复原每一条航程。读者从图回到文本时，应带着一个较小却
更可靠的问题：哪一种消息和物资可以走到哪里，谁使它继续走，何处又使它停下。
